\documentclass[11pt,a4paper]{article}

\usepackage[utf8]{inputenc}
\usepackage[T1]{fontenc}
\usepackage[french]{babel}
\usepackage[a4paper,margin=2.2cm]{geometry}
\usepackage{amsmath,amssymb,amsthm,mathtools}
\usepackage{lmodern}
\usepackage{enumitem}
\usepackage{hyperref}
\usepackage{fancyhdr}
\usepackage{xcolor}

\hypersetup{
    colorlinks=true,
    linkcolor=black,
    urlcolor=blue
}

\setlength{\parindent}{0pt}
\setlength{\parskip}{0.6em}

\pagestyle{fancy}
\fancyhf{}
\fancyhead[L]{Corrigé --- Nombres complexes}
\fancyhead[R]{Terminale $\to$ 1re année de prépa}
\fancyfoot[C]{\thepage}

\newcommand{\R}{\mathbb{R}}
\newcommand{\C}{\mathbb{C}}

\begin{document}

\begin{center}
    {\LARGE \textbf{Corrigé complet --- Nombres complexes}}\\[0.3em]
    {\large De la terminale à la première année de prépa}
\end{center}

\hrule
\vspace{0.8em}

\tableofcontents

\newpage

\section*{I. Forme algébrique et calculs}
\addcontentsline{toc}{section}{I. Forme algébrique et calculs}

\textbf{Exercice 1}

\[
z_1=3+2i \Rightarrow \Re(z_1)=3,\ \Im(z_1)=2.
\]
\[
z_2=-5i \Rightarrow \Re(z_2)=0,\ \Im(z_2)=-5.
\]
\[
z_3=7 \Rightarrow \Re(z_3)=7,\ \Im(z_3)=0.
\]
\[
z_4=-2+4i \Rightarrow \Re(z_4)=-2,\ \Im(z_4)=4.
\]

\textbf{Exercice 2}

\begin{itemize}
    \item \(2\) est réel.
    \item \(-3i\) est imaginaire pur.
    \item \(4+0i\) est réel.
    \item \(i\sqrt{5}\) est imaginaire pur.
    \item \(1-i\) n’est ni réel ni imaginaire pur.
\end{itemize}

\textbf{Exercice 3}

\begin{enumerate}[label=\alph*)]
    \item
    \[
    (2+3i)+(4-5i)=6-2i.
    \]

    \item
    \[
    (1-2i)(3+i)=3+i-6i-2i^2=3-5i+2=5-5i.
    \]

    \item
    \[
    (2+i)^2=4+4i+i^2=3+4i.
    \]

    \item
    \[
    (1+i)^2=1+2i+i^2=2i,
    \quad \text{donc} \quad
    (1+i)^3=(1+i)\cdot 2i=2i+2i^2=-2+2i.
    \]
\end{enumerate}

\textbf{Exercice 4}

\begin{enumerate}[label=\alph*)]
    \item
    \[
    \frac{1}{2+i}=\frac{2-i}{(2+i)(2-i)}=\frac{2-i}{5}.
    \]

    \item
    \[
    \frac{3-2i}{1+i}
    =\frac{(3-2i)(1-i)}{(1+i)(1-i)}
    =\frac{3-3i-2i+2i^2}{2}
    =\frac{1-5i}{2}.
    \]

    \item
    \[
    \frac{2+3i}{1-2i}
    =\frac{(2+3i)(1+2i)}{(1-2i)(1+2i)}
    =\frac{2+4i+3i+6i^2}{5}
    =\frac{-4+7i}{5}.
    \]
\end{enumerate}

\textbf{Exercice 5}

\[
z+2i=3-5i \Longrightarrow z=3-7i.
\]

\[
(1-i)z=4+2i \Longrightarrow z=\frac{4+2i}{1-i}
=\frac{(4+2i)(1+i)}{2}
=\frac{2+6i}{1}=1+3i? 
\]

Correction :
\[
(4+2i)(1+i)=4+4i+2i+2i^2=2+6i.
\]
Donc
\[
z=\frac{2+6i}{2}=1+3i.
\]

\textbf{Exercice 6}

Si
\[
z_1=a+ib,\qquad z_2=c+id,
\]
alors
\[
z_1z_2=(a+ib)(c+id)=ac+i(ad+bc)+i^2bd=(ac-bd)+i(ad+bc).
\]
Donc
\[
\Re(z_1z_2)=ac-bd=\Re(z_1)\Re(z_2)-\Im(z_1)\Im(z_2),
\]
\[
\Im(z_1z_2)=ad+bc=\Re(z_1)\Im(z_2)+\Im(z_1)\Re(z_2).
\]

\section*{II. Conjugué, module et propriétés}
\addcontentsline{toc}{section}{II. Conjugué, module et propriétés}

\textbf{Exercice 7}

\[
\overline{3+2i}=3-2i,\qquad \overline{-1+4i}=-1-4i,
\qquad \overline{5}=5,\qquad \overline{-7i}=7i.
\]

\textbf{Exercice 8}

Si \(z=a+ib\), alors \(\overline z=a-ib\).

\begin{enumerate}[label=\alph*)]
    \item
    \[
    z=\overline z
    \Longleftrightarrow a+ib=a-ib
    \Longleftrightarrow 2ib=0
    \Longleftrightarrow b=0.
    \]
    Donc \(z\in\R\).

    \item
    \[
    z=-\overline z
    \Longleftrightarrow a+ib=-a+ib
    \Longleftrightarrow 2a=0
    \Longleftrightarrow a=0.
    \]
    Donc \(z\in i\R\).
\end{enumerate}

\textbf{Exercice 9}

\[
|1+i|=\sqrt{1^2+1^2}=\sqrt{2},
\qquad
|-3+4i|=\sqrt{9+16}=5,
\]
\[
|2|=2,
\qquad
|-i\sqrt3|=\sqrt3.
\]

\textbf{Exercice 10}

Pour \(z=1+i\),
\[
z\overline z=(1+i)(1-i)=2=|z|^2.
\]

Pour \(z=-3+4i\),
\[
z\overline z=(-3+4i)(-3-4i)=9+16=25=|z|^2.
\]

\textbf{Exercice 11}

Si \(z=a+ib\), alors
\[
|z|=\sqrt{a^2+b^2}.
\]
Comme
\[
a^2\le a^2+b^2 \quad \text{et} \quad b^2\le a^2+b^2,
\]
on obtient
\[
|a|\le \sqrt{a^2+b^2}=|z|,
\qquad
|b|\le \sqrt{a^2+b^2}=|z|.
\]
Donc
\[
|\Re(z)|\le |z|,
\qquad
|\Im(z)|\le |z|.
\]

\textbf{Exercice 12}

Si \(z\neq 0\), alors
\[
z\overline z=|z|^2.
\]
Donc
\[
z\cdot \frac{\overline z}{|z|^2}=1,
\]
ce qui prouve que
\[
\frac{1}{z}=\frac{\overline z}{|z|^2}.
\]

\textbf{Exercice 13}

On calcule :
\[
|zz'|^2=(zz')(\overline{zz'})=zz'\overline z\,\overline{z'}
=(z\overline z)(z'\overline{z'})=|z|^2|z'|^2.
\]
Les deux membres étant positifs :
\[
|zz'|=|z||z'|.
\]

Si \(z'\neq 0\),
\[
\left|\frac{z}{z'}\right|
=\left|z\cdot \frac{1}{z'}\right|
=|z|\cdot \left|\frac{1}{z'}\right|
=|z|\cdot \frac{1}{|z'|}
=\frac{|z|}{|z'|}.
\]

\section*{III. Interprétation géométrique}
\addcontentsline{toc}{section}{III. Interprétation géométrique}

\textbf{Exercice 14}

\[
A(2,3)\longleftrightarrow z_A=2+3i,
\qquad
B(-1,4)\longleftrightarrow z_B=-1+4i,
\qquad
C(0,-2)\longleftrightarrow z_C=-2i.
\]

\textbf{Exercice 15}

\[
z_{\overrightarrow{AB}}=z_B-z_A=(4-i)-(1+2i)=3-3i.
\]

\textbf{Exercice 16}

\begin{enumerate}[label=\alph*)]
    \item
    \[
    z_I=\frac{z_A+z_B}{2}
    =\frac{(2+i)+(4-3i)}{2}
    =\frac{6-2i}{2}=3-i.
    \]

    \item
    \[
    z_I=\frac{(-1+5i)+(3+i)}{2}
    =\frac{2+6i}{2}=1+3i.
    \]
\end{enumerate}

\textbf{Exercice 17}

Le vecteur \(\overrightarrow{AB}\) a pour affixe
\[
z_B-z_A.
\]
Sa norme vaut donc
\[
|z_B-z_A|.
\]
Or la longueur \(AB\) est la norme du vecteur \(\overrightarrow{AB}\). Donc
\[
AB=|z_B-z_A|.
\]

\section*{IV. Inégalité triangulaire}
\addcontentsline{toc}{section}{IV. Inégalité triangulaire}

\textbf{Exercice 18}

\[
z+z'=(1+i)+(1-i)=2,
\quad \text{donc} \quad |z+z'|=2.
\]
\[
|z|=\sqrt2,\qquad |z'|=\sqrt2,
\quad \text{donc} \quad |z|+|z'|=2\sqrt2.
\]
On a bien
\[
|z+z'|\le |z|+|z'|.
\]

\textbf{Exercice 19}

On calcule :
\[
|z+z'|^2=(z+z')(\overline z+\overline{z'})
=|z|^2+|z'|^2+z\overline{z'}+\overline z z'.
\]
Or
\[
z\overline{z'}+\overline z z'=2\Re(z\overline{z'}),
\]
donc
\[
|z+z'|^2=|z|^2+|z'|^2+2\Re(z\overline{z'}).
\]
Comme
\[
\Re(w)\le |w|,
\]
on obtient
\[
|z+z'|^2\le |z|^2+|z'|^2+2|z\overline{z'}|.
\]
Mais
\[
|z\overline{z'}|=|z||z'|.
\]
Ainsi
\[
|z+z'|^2\le |z|^2+|z'|^2+2|z||z'|=(|z|+|z'|)^2.
\]
Comme les modules sont positifs,
\[
|z+z'|\le |z|+|z'|.
\]

\textbf{Exercice 20}

On écrit
\[
z=(z-z')+z'.
\]
Par l’inégalité triangulaire,
\[
|z|\le |z-z'|+|z'|,
\]
donc
\[
|z|-|z'|\le |z-z'|.
\]
En échangeant \(z\) et \(z'\), on obtient
\[
|z'|-|z|\le |z-z'|.
\]
D’où
\[
\big||z|-|z'|\big|\le |z-z'|.
\]

\section*{V. Forme exponentielle et trigonométrique}
\addcontentsline{toc}{section}{V. Forme exponentielle et trigonométrique}

\textbf{Exercice 21}

\[
1=e^{i0},
\qquad
-1=e^{i\pi},
\qquad
i=e^{i\pi/2},
\qquad
-i=e^{-i\pi/2}.
\]

\textbf{Exercice 22}

\begin{enumerate}[label=\alph*)]
    \item
    \[
    z_1=1+i,\quad |z_1|=\sqrt2,\quad \arg(z_1)=\frac{\pi}{4}\ [2\pi].
    \]

    \item
    \[
    z_2=-1+i,\quad |z_2|=\sqrt2,\quad \arg(z_2)=\frac{3\pi}{4}\ [2\pi].
    \]

    \item
    \[
    z_3=-\sqrt3+i,\quad |z_3|=2.
    \]
    On a
    \[
    \cos\theta=-\frac{\sqrt3}{2},\qquad \sin\theta=\frac12,
    \]
    donc
    \[
    \arg(z_3)=\frac{5\pi}{6}\ [2\pi].
    \]

    \item
    \[
    z_4=2e^{i\pi/3},\quad |z_4|=2,\quad \arg(z_4)=\frac{\pi}{3}\ [2\pi].
    \]
\end{enumerate}

\textbf{Exercice 23}

\[
z=-1+i\sqrt3.
\]
Son module vaut
\[
|z|=\sqrt{1+3}=2.
\]
On cherche \(\theta\) tel que
\[
\cos\theta=-\frac12,\qquad \sin\theta=\frac{\sqrt3}{2}.
\]
Donc
\[
\theta=\frac{2\pi}{3}\ [2\pi].
\]
Ainsi
\[
z=2\left(\cos\frac{2\pi}{3}+i\sin\frac{2\pi}{3}\right)=2e^{i2\pi/3}.
\]

\textbf{Exercice 24}

\[
e^{ia}e^{ib}=(\cos a+i\sin a)(\cos b+i\sin b).
\]
En développant :
\[
e^{ia}e^{ib}=(\cos a\cos b-\sin a\sin b)+i(\sin a\cos b+\cos a\sin b).
\]
Par les formules d’addition,
\[
e^{ia}e^{ib}=\cos(a+b)+i\sin(a+b)=e^{i(a+b)}.
\]

\textbf{Exercice 25}

On a
\[
e^{i\theta}=\cos\theta+i\sin\theta,
\qquad
e^{-i\theta}=\cos\theta-i\sin\theta.
\]
En additionnant :
\[
e^{i\theta}+e^{-i\theta}=2\cos\theta
\Longrightarrow
\cos\theta=\frac{e^{i\theta}+e^{-i\theta}}{2}.
\]
En soustrayant :
\[
e^{i\theta}-e^{-i\theta}=2i\sin\theta
\Longrightarrow
\sin\theta=\frac{e^{i\theta}-e^{-i\theta}}{2i}.
\]

\textbf{Exercice 26}

\begin{enumerate}[label=\alph*)]
    \item Par Moivre,
    \[
    (\cos\theta+i\sin\theta)^4=\cos(4\theta)+i\sin(4\theta).
    \]

    \item
    \[
    1+i=\sqrt2\,e^{i\pi/4},
    \]
    donc
    \[
    (1+i)^5=(\sqrt2)^5 e^{i5\pi/4}=4\sqrt2\,e^{i5\pi/4}.
    \]
    Sous forme algébrique :
    \[
    (1+i)^5=4\sqrt2\left(-\frac{\sqrt2}{2}-i\frac{\sqrt2}{2}\right)=-4-4i.
    \]
\end{enumerate}

\section*{VI. Équations, racines carrées et racines \(n\)-ièmes}
\addcontentsline{toc}{section}{VI. Équations, racines carrées et racines n-ièmes}

\textbf{Exercice 27}

\[
z^2=-4 \Longrightarrow z=\pm 2i.
\]
\[
z^2=-9 \Longrightarrow z=\pm 3i.
\]
\[
z^2=-16 \Longrightarrow z=\pm 4i.
\]

\textbf{Exercice 28}

\begin{enumerate}
    \item
    \[
    z^2+1=0 \Longrightarrow z^2=-1 \Longrightarrow z=\pm i.
    \]

    \item
    \[
    z^2-2z+5=0.
    \]
    Discriminant :
    \[
    \Delta=(-2)^2-4\cdot 1\cdot 5=4-20=-16.
    \]
    Donc
    \[
    \sqrt{\Delta}=\pm 4i.
    \]
    Les solutions sont
    \[
    z=\frac{2\pm 4i}{2}=1\pm 2i.
    \]
\end{enumerate}

\textbf{Exercice 29}

On résout
\[
z^2-(1+4i)z+(3i-3)=0.
\]
Ici
\[
a=1,\quad b=-(1+4i),\quad c=3i-3.
\]
Le discriminant vaut
\[
\Delta=b^2-4ac=( -1-4i)^2-4(3i-3).
\]
Or
\[
(-1-4i)^2=1+8i+16i^2=-15+8i.
\]
Donc
\[
\Delta=(-15+8i)-12i+12=-3-4i.
\]
On cherche \(\delta\) tel que \(\delta^2=-3-4i\). On reconnaît
\[
(1-2i)^2=1-4i-4=-3-4i.
\]
Donc \(\delta=\pm(1-2i)\).

Les solutions sont alors
\[
z=\frac{1+4i\pm (1-2i)}{2}.
\]
Ainsi :
\[
z_1=\frac{2+2i}{2}=1+i,
\qquad
z_2=\frac{6i}{2}=3i.
\]

\textbf{Exercice 30}

Les racines cubiques de l’unité sont les solutions de
\[
z^3=1=e^{i0}.
\]
Donc
\[
z_k=e^{2ik\pi/3},\qquad k=0,1,2.
\]
Ainsi :
\[
1,\qquad e^{2i\pi/3}=-\frac12+i\frac{\sqrt3}{2},\qquad
e^{4i\pi/3}=-\frac12-i\frac{\sqrt3}{2}.
\]

\textbf{Exercice 31}

On écrit
\[
-16=16e^{i\pi}.
\]
Les racines quatrièmes sont
\[
z_k=\sqrt[4]{16}\,e^{i(\pi+2k\pi)/4}=2e^{i(\pi/4+k\pi/2)},
\qquad k=0,1,2,3.
\]
Donc
\[
2e^{i\pi/4},\quad 2e^{i3\pi/4},\quad 2e^{i5\pi/4},\quad 2e^{i7\pi/4}.
\]

\textbf{Exercice 32}

Si
\[
Z=\rho e^{i\theta}\quad (\rho>0),
\]
alors les solutions de
\[
z^n=Z
\]
sont
\[
z_k=\sqrt[n]{\rho}\,e^{i(\theta+2k\pi)/n},\qquad k=0,1,\dots,n-1.
\]

\section*{VII. Transformations trigonométriques}
\addcontentsline{toc}{section}{VII. Transformations trigonométriques}

\textbf{Exercice 33}

On veut écrire
\[
2\cos t+\sqrt3\sin t=R\cos(t-\varphi).
\]
Or
\[
R\cos(t-\varphi)=R(\cos t\cos\varphi+\sin t\sin\varphi).
\]
On identifie :
\[
R\cos\varphi=2,\qquad R\sin\varphi=\sqrt3.
\]
Donc
\[
R=\sqrt{4+3}=\sqrt7.
\]
Puis
\[
\cos\varphi=\frac{2}{\sqrt7},\qquad \sin\varphi=\frac{\sqrt3}{\sqrt7}.
\]
Ainsi
\[
2\cos t+\sqrt3\sin t=\sqrt7\,\cos(t-\varphi),
\]
avec
\[
\varphi=\arctan\left(\frac{\sqrt3}{2}\right)
\]
ou, plus rigoureusement, \(\varphi\) défini par les deux relations ci-dessus.

\textbf{Exercice 34}

On écrit
\[
-\cos t+2\sin t=R\cos(t-\varphi).
\]
En identifiant :
\[
R\cos\varphi=-1,\qquad R\sin\varphi=2.
\]
Donc
\[
R=\sqrt{1+4}=\sqrt5.
\]
On choisit \(\varphi\) tel que
\[
\cos\varphi=-\frac{1}{\sqrt5},\qquad \sin\varphi=\frac{2}{\sqrt5}.
\]
Donc
\[
-\cos t+2\sin t=\sqrt5\,\cos(t-\varphi).
\]

On peut aussi écrire
\[
-\cos t+2\sin t=R\sin(t+\psi)
=R(\sin t\cos\psi+\cos t\sin\psi).
\]
On identifie :
\[
R\cos\psi=2,\qquad R\sin\psi=-1.
\]
Donc
\[
R=\sqrt5,
\qquad
\cos\psi=\frac{2}{\sqrt5},\qquad \sin\psi=-\frac{1}{\sqrt5}.
\]
Ainsi
\[
-\cos t+2\sin t=\sqrt5\,\sin(t+\psi).
\]

\section*{VIII. Exercices de synthèse}
\addcontentsline{toc}{section}{VIII. Exercices de synthèse}

\textbf{Exercice 35}

Si \(|z|=1\), alors
\[
z\overline z=|z|^2=1.
\]
Donc \(\overline z\) est l’inverse de \(z\), d’où
\[
\overline z=\frac{1}{z}.
\]

\textbf{Exercice 36}

On a
\[
\left|\frac{z}{\overline z}\right|=\frac{|z|}{|\overline z|}.
\]
Or
\[
|\overline z|=|z|.
\]
Donc
\[
\left|\frac{z}{\overline z}\right|=1.
\]

\textbf{Exercice 37}

Posons
\[
z=x+iy.
\]
Alors
\[
z+\overline z=(x+iy)+(x-iy)=2x.
\]
La condition
\[
z+\overline z=4
\]
donne
\[
2x=4 \Longrightarrow x=2.
\]
De plus
\[
|z|=3 \Longrightarrow x^2+y^2=9.
\]
Donc
\[
4+y^2=9 \Longrightarrow y^2=5 \Longrightarrow y=\pm \sqrt5.
\]
Les solutions sont donc
\[
z=2+i\sqrt5
\qquad \text{ou} \qquad
z=2-i\sqrt5.
\]

\textbf{Exercice 38}

Résoudre
\[
z^4=1.
\]
Les solutions sont les racines quatrièmes de l’unité :
\[
z_k=e^{2ik\pi/4}=e^{ik\pi/2},\qquad k=0,1,2,3.
\]
Donc
\[
1,\quad i,\quad -1,\quad -i.
\]
Elles sont situées sur le cercle unité aux angles
\[
0,\ \frac{\pi}{2},\ \pi,\ \frac{3\pi}{2}.
\]

\textbf{Exercice 39}

On écrit
\[
-64=64e^{i(\pi+2m\pi)}.
\]
Les solutions de
\[
z^6=-64
\]
sont donc
\[
z_k=\sqrt[6]{64}\,e^{i(\pi+2k\pi)/6}=2e^{i(\pi+2k\pi)/6},
\qquad k=0,1,2,3,4,5.
\]
Autrement dit
\[
z_k=2e^{i(\pi/6+k\pi/3)},\qquad k=0,\dots,5.
\]

\textbf{Exercice 40}

Les racines \(n\)-ièmes de l’unité sont
\[
\omega_k=e^{2ik\pi/n},\qquad k=0,1,\dots,n-1.
\]
Elles ont toutes pour module \(1\), donc elles appartiennent au cercle unité.

Leurs arguments sont régulièrement espacés :
\[
0,\ \frac{2\pi}{n},\ \frac{4\pi}{n},\ \dots,\ \frac{2(n-1)\pi}{n}.
\]
L’écart angulaire entre deux racines consécutives est constant et vaut
\[
\frac{2\pi}{n}.
\]
Elles sont donc les sommets d’un polygone régulier à \(n\) côtés inscrit dans le cercle unité.

\end{document}